\chapter{System Design and Architecture}
\label{chap:design}

\section{Design Philosophy}
\label{sec:design-philosophy}

Before getting into the specifics of each subsystem, it is worth saying something
about what we were trying to avoid. The most common failure mode in student AI
projects is what might be called ``glue-code'' architecture: take an LLM API,
wrap a thin web layer around it, and call it a platform. That approach produces
something that technically works but breaks down immediately under any real-world
load, is impossible to diagnose when it produces bad outputs, and offers no
foundation for improvement.

We deliberately went in the other direction. Every layer of EngUnity has a reason
for being where it is and a defined boundary with its neighbours. When a generated
answer is poor, we can point to a specific module---query classifier, re-ranker,
compression step, refiner---and ask what went wrong there. That traceability was
a deliberate design goal, not an afterthought.

\section{Overall Architecture}
\label{sec:architecture}

EngUnity is organised into four planes that communicate only through typed API or
message-queue contracts. Keeping these boundaries strict is what makes individual
components replaceable---if FAISS becomes a bottleneck, swapping in Qdrant should
not require changes beyond the vector-store interface.

\begin{table}[ht]
    \centering
    \caption{The four architectural planes and their primary responsibilities.}
    \label{tab:planes}
    \begin{tabular}{lll}
        \toprule
        \textbf{Plane} & \textbf{Technology} & \textbf{Responsibilities} \\
        \midrule
        Presentation   & Next.js 16, React 18, Tailwind CSS 3 &
            Rendering, routing, real-time SSE consumption \\
        Orchestration  & FastAPI 0.115, Uvicorn, Celery, Redis &
            Auth, rate-limiting, job dispatch, API contracts \\
        Inference      & OmniRAG pipeline, AI agents, Groq / Ollama &
            Retrieval, reasoning, generation, quality evaluation \\
        Persistence    & PostgreSQL, MongoDB, Redis, FAISS &
            Relational data, document store, cache, vector index \\
        \bottomrule
    \end{tabular}
\end{table}

The client (Next.js PWA) communicates exclusively with the Orchestration Plane
over HTTPS. The Inference and Persistence planes are never directly accessible
from the browser. Background jobs---file indexing, bulk analysis, knowledge-graph
community building---go through Celery, backed by Redis, so that long-running
operations do not block the FastAPI event loop.

\section{The OmniRAG Pipeline}
\label{sec:omnirag}

\subsection{Why a Single Pipeline for Four Strategies?}

Early in the project we had separate endpoints for simple questions, document
search, and deep research. This worked, but it pushed a significant cognitive
burden onto the student: they had to decide what \textit{kind} of question they
were asking before they could get a useful answer. A casual empirical check
showed that students routinely underestimated query complexity and ended up on
the wrong endpoint. The fix was to classify automatically and route internally.
The pipeline exposes a single \texttt{process\_query} method; the strategy is
an implementation detail.

\subsection{Complexity Classification and Routing}

The \texttt{QueryComplexityClassifier} analyses each incoming query and assigns
one of four classes:

\begin{table}[ht]
    \centering
    \caption{OmniRAG complexity classes, strategies, and typical query characteristics.}
    \label{tab:complexity}
    \begin{tabular}{llp{6cm}}
        \toprule
        \textbf{Class} & \textbf{Strategy} & \textbf{Characteristics} \\
        \midrule
        \texttt{SIMPLE} & Direct generation &
            Conversational, definitional, low document-dependency \\
        \texttt{SINGLE\_HOP} & Vector RAG &
            Factual with a clear document anchor; single retrieval topic \\
        \texttt{MULTI\_HOP} & Graph RAG (map-reduce) &
            Comparative, synthesising, requires multiple source documents \\
        \texttt{RECURSIVE\_INTENSIVE} & Recursive reasoning &
            Open-ended analytical, benefits from multi-step hypothesis chains \\
        \bottomrule
    \end{tabular}
\end{table}

\subsection{The Single-Hop Retrieval Flow}

For the most common class of queries---single-hop---the pipeline executes the
following steps. We have numbered them to make it easier to cross-reference with
the source code in \texttt{backend/app/services/rag/pipeline.py}.

\begin{enumerate}
    \item \textbf{Memory recall.} Up to three relevant memories are fetched from
    the user's long-term profile via \texttt{memory\_system.recall}. These are
    injected into the query rewriting step.

    \item \textbf{Image perception (conditional).} If the request includes
    \texttt{image\_urls} or \texttt{image\_ids}, the pipeline calls the
    \texttt{MultiModalVisionProcessor} before any text retrieval occurs. The
    processor returns a structured visual context string that is prepended to
    the query.

    \item \textbf{Query rewriting.} The \texttt{QueryRewriter} reformulates
    the raw user query---incorporating the memory context from step 1---to
    improve retrieval recall. This step is particularly useful for follow-up
    questions that omit pronouns or assume prior context.

    \item \textbf{Multi-query expansion.} The LLM generates four query
    variations including a step-back abstraction (a broader, more general form
    of the original question). Retrieval runs concurrently for each variant
    using \texttt{asyncio.gather}.

    \item \textbf{HyDE transformation.} \texttt{HyDEEngine} generates a
    hypothetical answer and embeds it. The resulting embedding is used as the
    retrieval vector alongside the original query embedding. In our testing on
    technical engineering questions, HyDE consistently retrieved more relevant
    passage excerpts than the raw query embedding alone.

    \item \textbf{Hybrid search.} FAISS performs a combined semantic and BM25
    search with a weighting of $\alpha=0.6$ semantic and $0.4$ keyword. Top-20
    results are retrieved per query variant and then deduplicated by chunk ID.

    \item \textbf{FlashRank neural re-ranking.} The cross-attention re-ranker
    in \texttt{FlashRankReranker} scores all deduplicated candidates and
    retains the top 10. This step is what separates passages that are
    topically related from those that are genuinely answer-relevant.

    \item \textbf{CRAG evaluation.} \texttt{CRAGPipeline} scores retrieval
    quality using the \texttt{RetrievalEvaluator}. If the top document scores
    below the configured threshold, \texttt{WebSearchFallback} queries the
    Brave Search API and supplements the candidate pool.

    \item \textbf{Contextual compression.} Each of the top-5 re-ranked
    documents is compressed by an LLM call that extracts only the excerpt
    most relevant to the query. The five compression calls run in parallel;
    this step typically eliminates 60--70\% of each document's token count
    while preserving the answer-bearing sentences.

    \item \textbf{Two-stage generation.} A draft is generated at
    \texttt{temperature=0.3} using the compressed context. The
    \texttt{AnswerRefiner} then inspects the draft and removes filler phrases,
    reduces passive constructions, and trims word count while preserving
    factual content.

    \item \textbf{Self-critique.} \texttt{SelfCritique} scores the refined
    answer against the retrieved documents and emits a confidence value.
    Answers that score poorly here are flagged in the response metadata so
    that the frontend can surface a visual indicator to the user.

    \item \textbf{Quality metric logging.} Structure, density, naturalness,
    and confidence scores are computed and appended to
    \texttt{quality\_metrics.jsonl} for offline analysis.
\end{enumerate}

\subsection{Multi-Hop (Graph RAG) Flow}

The multi-hop path retrieves knowledge-graph community summaries alongside
standard vector chunks and synthesises them using a map-reduce pattern.

In the \textit{map} phase, the pipeline generates partial answers independently
for each community summary and each top-5 document, running all LLM calls in
parallel via \texttt{asyncio.gather}. In the \textit{reduce} phase, a
synthesis prompt receives all partial answers and produces a single, coherent
final response. The map-reduce structure means that the number of parallel LLM
calls scales with the number of retrieved communities and documents, which is
configurable at deployment time.

\subsection{Recursive Intensive Flow}

Recursive queries expand the retrieval budget ($k=40$ versus $k=20$ for single-
hop) and combine the larger document pool with knowledge-graph context. The
assembled context is handed to \texttt{RecursiveReasoningAgent}, which implements
a chain-of-thought loop where each step can spawn subsidiary inferences before
committing to a final answer. Each intermediate step is emitted as a structured
streaming event, giving the user partial visibility into the reasoning process.

\section{The DeepResearchAgent}
\label{sec:deepresearch}

\subsection{Design Rationale}

The OmniRAG pipeline, even on its recursive-intensive path, operates within a
single-call budget. For genuine research tasks---``summarise the state of the art
in thermophotovoltaic energy conversion'' or ``compare the memory management
approaches of five major programming languages''---this is not enough. The research
agent is the answer to that limitation.

The core design principle is simple: a single research iteration is like any other
single-hop retrieval, but the agent can run multiple iterations, each informed by
what the previous iterations found (and, more importantly, what they did not find).

\subsection{The Five-Phase Pipeline}

\begin{enumerate}
    \item \textbf{Decompose.} The main query is broken down by the LLM into
    up to $N$ typed sub-questions (factual, analytical, comparative, exploratory).
    The typing matters: factual sub-questions are sent to the vector store first;
    analytical and comparative ones are weighted toward the knowledge graph.

    \item \textbf{Search.} For each sub-question, three backends are queried
    concurrently: the OmniRAG vector store, a live web search (Brave Search API),
    and the knowledge-graph community index. Results from all three are pooled.

    \item \textbf{Evaluate.} A hybrid relevance scorer applies fast keyword-overlap
    matching as a pre-filter, then invokes an LLM scorer for ambiguous candidates.
    Snippets below a relevance threshold of 0.4 are discarded. The threshold is
    conservative because false positives at this stage degrade synthesis quality.

    \item \textbf{Refine.} The agent inspects which sub-questions remain
    under-sourced (fewer than the configured minimum source count). For each gap,
    it generates one or more exploratory follow-up queries and re-enters the
    Search phase. If the iteration budget is exhausted without reaching minimum
    coverage, the agent continues with whatever sources it has.

    \item \textbf{Synthesize.} All accepted snippets are assembled into a structured
    report prompt that produces: Executive Summary, Key Findings, Analysis,
    Caveats, Follow-up Questions, and Related Topics. The Caveats section is
    produced because early tests found that leaving known limitations implicit
    generated less trust from students than surfacing them explicitly.
\end{enumerate}

\begin{table}[ht]
    \centering
    \caption{ResearchDepth parameter matrix.}
    \label{tab:depth}
    \begin{tabular}{lrrrl}
        \toprule
        \textbf{Depth} & \textbf{Max Sub-Q} & \textbf{Max Iter} &
        \textbf{Min Src} & \textbf{Typical Use} \\
        \midrule
        \texttt{QUICK}      & 2  & 1 & 2  & Quick factual look-up \\
        \texttt{STANDARD}   & 5  & 3 & 4  & Coursework research \\
        \texttt{DEEP}       & 8  & 5 & 7  & Thesis-level questions \\
        \texttt{EXHAUSTIVE} & 12 & 8 & 10 & Full domain survey \\
        \bottomrule
    \end{tabular}
\end{table}

\section{Multimodal Visual Perception}
\label{sec:multimodal}

The \texttt{MultiModalVisionProcessor} wraps three perception components:

\begin{itemize}
    \item \textbf{CLIP (ViT-B/32)} embeds each image into the same 512-dimensional
    space used for text embeddings, enabling cross-modal retrieval within the same
    FAISS index.

    \item \textbf{EasyOCR} extracts printed and handwritten text with layout
    information preserved. Detected text regions are concatenated with positional
    markers (top-left, bottom-right coordinates) so that spatial queries remain
    meaningful.

    \item \textbf{YOLOv8n} detects objects and classifies them using the 80-class
    COCO taxonomy. Detected labels are included in the visual context string;
    for engineering images, the most useful classes in practice have been
    ``keyboard'', ``monitor'', ``book'', ``cell phone'', and various electronic
    components that fall back to ``remote'' or ``cell phone'' under the COCO
    taxonomy (a known limitation---see Chapter~\ref{chap:results}).
\end{itemize}

All three components produce text output that is concatenated into a visual context
prefix injected into the RAG prompt. This architecture means the same pipeline
handles text-only and multimodal queries without branching.

\section{Code Laboratory}
\label{sec:codelab}

The Code Lab is built around Monaco Editor (v0.50) on the frontend and a
\texttt{SecureSandbox} class on the backend. The sandbox spins up a Docker
container per execution request, enforces CPU and memory cgroups, and redirects
stdin and stdout through a controlled channel. There is no network access from
within the container by default.

On top of the execution layer:

\begin{itemize}
    \item A WebSocket endpoint at \texttt{/ws/terminal/\{project\_id\}} provides
    an interactive terminal through XTerm.js. The same container that ran the
    initial code execution remains active for the terminal session.

    \item AI inline completion streams token-by-token code suggestions from the
    OmniRAG pipeline and displays them as ghost text in Monaco. The completion
    prompt includes the current file content, cursor position, and a summary
    of recently visited files for context.

    \item A Debug Adapter Protocol implementation using Python's \texttt{Bdb}
    supports breakpoints, step-over, step-into, and variable inspection for
    Python programs. The DAP protocol messages are transported over the
    existing WebSocket connection.

    \item Git integration (\texttt{/api/v1/git/}) exposes log, diff, commit,
    and branch operations backed by GitPython. GitHub synchronisation through
    PyGithub 2.8 allows students to pull in remote repositories for analysis
    or push work back to their organisation.
\end{itemize}

\section{Memory Architecture}
\label{sec:memory}

Two memory layers operate in parallel:

\textbf{Short-term session summaries.} At natural pause points within a session,
the LLM generates a hierarchical summary of the conversation history. This
summary---rather than the raw message log---is passed as \texttt{memory\_summary}
on subsequent requests, keeping the context window lean while preserving the
substance of prior exchanges.

\textbf{Long-term user profiles.} After each OmniRAG interaction, the
\texttt{memory\_system.remember()} call stores a structured record: the original
query, the generated response, the sources cited, the complexity class and
strategy used, and the self-critique confidence score. A background aggregation
process computes per-user patterns from this record store and populates a profile
that the query rewriter consults on subsequent sessions.

Both layers are scoped to individual users by \texttt{user\_id}. There is no
cross-user memory sharing.

\section{Analytics Workbench}
\label{sec:analytics}

The Analytics module accepts CSV and XLSX uploads, infers column types using
Pandas, and provides a self-service workbench for descriptive statistics,
K-Means clustering (scikit-learn), and lightweight classification. Recharts
renders the output charts in the browser. The AI insight generation endpoint
feeds a dataset summary into the OmniRAG pipeline --- the same pipeline that
handles all other queries --- and returns a structured analysis covering
correlations, outliers, and recommended follow-up investigations. Export to
PDF is handled by jsPDF with jspdf-autotable.

\section{Decision Vault}
\label{sec:decision}

The Decision Vault stores structured decisions (context, constraints, options,
rationale, confidence) and submits each to \texttt{DecisionAIService} for
adversarial review. The AI reviewer attempts to challenge the premises of the
decision and flag logical inconsistencies or hidden assumptions. Results are
returned with severity labels (critical, warning, suggestion) and are
exportable as PDF or JSON.

This subsystem was the least mature of the five at the time of writing. The
review logic works, but the prompt engineering for adversarial critiquing is an
area where iteration is clearly warranted.

\section{Authentication and Security}
\label{sec:security}

Authentication uses two cooperating layers. Supabase handles the primary auth
flow (email/password and OAuth2 social providers); every request to the FastAPI
backend carries a Supabase-issued JWT that is validated by
\texttt{get\_current\_user()}. For relational integrity, each successful Supabase
authentication also upserts a local user record in PostgreSQL, so that
foreign-key relationships (sessions, projects, analysis records) do not depend
on a remote service being available at query time.

Security testing covers input validation (injection attacks on all auth
endpoints via \texttt{TestInputValidation}), standard header assertions
(\texttt{TestSecurityHeaders}), JWT-specific tests (\texttt{TestJWTTokens}),
and password-hash correctness (\texttt{TestPasswordHashing}). Rate limiting
is enforced globally via SlowAPI backed by Redis, with per-user overrides
configurable through environment variables.
